% GENERATED FILE. Edit canonical lessons, then run npm run generate:books.
% canonical-source-hash: fnv1a64:6e37e9fa07701b4e
% canonical-lessons: FR-C311-fiancer, FR-C311-demande-en-mariage, FR-C311-marier, FR-C311-piece-montee, FR-C311-dragee

\chapter{Un grand jour}
\label{ch:fr-un-grand-jour}
\hlchaptermodality{311}

\begin{chapteropening}
\textbf{By the end of this chapter:} \emph{I can talk about getting engaged and married: getting engaged, the proposal, getting married, the wedding cake and sugared almonds.}

Complete the last lesson of chapter 311: I can talk about getting engaged and married: getting engaged, the proposal, getting married, the wedding cake and sugared almonds.
\end{chapteropening}

\section[fiancer]{fiancer — to get engaged (as se fiancer)}

\label{lesson:FR-C311-fiancer}



\begin{quote}
Before the new one: say the French for the start of the school year, then the French for of age.
\end{quote}

\subsection*{You'll want to know: fiancer}
\textbf{fiancer} — \textquotedblleft{}to get engaged (as se fiancer)\textquotedblright{}.

\begin{quote}
Ils se sont fiancés l'été dernier.
\end{quote}

\begin{quote}
\emph{They got engaged last summer.}
\end{quote}

Mostly reflexive: \emph{se fiancer avec quelqu'un}. The person is \emph{le fiancé}, \emph{la fiancée}.

\begin{cousinweb}
\emph{fiancer} is built on Old French \emph{fiance} \textquotedblleft{}trust, a promise\textquotedblright{}, from \emph{fier} \textquotedblleft{}to trust\textquotedblright{}, which goes back to Latin \emph{fīdus} \textquotedblleft{}faithful\textquotedblright{} (Lewis \& Short).
\end{cousinweb}

\begin{grammarlens}[title={What you've built}]
One more word for telling, explaining and describing at B1.
\end{grammarlens}

\subsection*{Your turn}
\begin{itemize}
\raggedright
  \item \emph{Say it:} \emph{fiancer}
  \item \emph{Say it:} \emph{fiancer}, once more
  \item \emph{Say it:} a sentence of your own with \emph{fiancer}
  \item \emph{Recall:} say the French for the start of the school year, then the French for of age, then say \emph{fiancer} again
\end{itemize}

\begin{tcolorbox}[breakable,colback=teal!4,colframe=teal!35!black,title={Before you move on}]
What does fiancer mean? (\textquotedblleft{}To get engaged (as se fiancer)\textquotedblright{}.) Say it once more.
\end{tcolorbox}
\section[la demande en mariage]{la demande en mariage — the marriage proposal}

\label{lesson:FR-C311-demande-en-mariage}



\begin{quote}
Before the new one: say the French for of age, then the French for to get engaged.
\end{quote}

\subsection*{You'll want to know: la demande en mariage}
\textbf{la demande en mariage} — \textquotedblleft{}the marriage proposal\textquotedblright{}.

\begin{quote}
Il a fait sa demande en mariage au restaurant.
\end{quote}

\begin{quote}
\emph{He proposed at the restaurant.}
\end{quote}

\emph{Demander quelqu'un en mariage} — to propose to someone.

\begin{cousinweb}
\emph{mariage} is built on \emph{marier}, inherited from Latin \emph{marītāre} \textquotedblleft{}to wed\textquotedblright{} (EtymDB); \emph{demande} is built on \emph{demander}, from Latin \emph{dēmandāre} \textquotedblleft{}to entrust\textquotedblright{} (Lewis \& Short).
\end{cousinweb}

\begin{grammarlens}[title={What you've built}]
One more word for telling, explaining and describing at B1.
\end{grammarlens}

\subsection*{Your turn}
\begin{itemize}
\raggedright
  \item \emph{Say it:} \emph{la demande en mariage}
  \item \emph{Say it:} \emph{la demande en mariage}, once more
  \item \emph{Say it:} a sentence of your own with \emph{la demande en mariage}
  \item \emph{Recall:} say the French for of age, then the French for to get engaged, then say \emph{la demande en mariage} again
\end{itemize}

\begin{tcolorbox}[breakable,colback=teal!4,colframe=teal!35!black,title={Before you move on}]
What does la demande en mariage mean? (\textquotedblleft{}The marriage proposal\textquotedblright{}.) Say it once more.
\end{tcolorbox}
\section[marier]{marier — to get married (as se marier); to marry (a couple), to marry off}

\label{lesson:FR-C311-marier}



\begin{quote}
Before the new one: say the French for to get engaged, then the French for the marriage proposal.
\end{quote}

\subsection*{You'll want to know: marier}
\textbf{marier} — \textquotedblleft{}to get married (as se marier); to marry (a couple), to marry off\textquotedblright{}.

\begin{quote}
Ils se sont mariés à la mairie.
\end{quote}

\begin{quote}
\emph{They got married at the town hall.}
\end{quote}

Reflexive for the couple: \emph{ils se marient en juin}. Without \emph{se}, the mayor marries them: \emph{le maire les a mariés}.

\begin{cousinweb}
\emph{marier} is inherited from Latin \emph{marītāre} \textquotedblleft{}to wed\textquotedblright{} (EtymDB), built on \emph{marītus} \textquotedblleft{}husband\textquotedblright{}.
\end{cousinweb}

\begin{grammarlens}[title={What you've built}]
One more word for telling, explaining and describing at B1.
\end{grammarlens}

\subsection*{Your turn}
\begin{itemize}
\raggedright
  \item \emph{Say it:} \emph{marier}
  \item \emph{Say it:} \emph{marier}, once more
  \item \emph{Say it:} a sentence of your own with \emph{marier}
  \item \emph{Recall:} say the French for to get engaged, then the French for the marriage proposal, then say \emph{marier} again
\end{itemize}

\begin{tcolorbox}[breakable,colback=teal!4,colframe=teal!35!black,title={Before you move on}]
What does marier mean? (\textquotedblleft{}To get married (as se marier); to marry (a couple), to marry off\textquotedblright{}.) Say it once more.
\end{tcolorbox}
\section[la pièce montée]{la pièce montée — the tiered wedding cake}

\label{lesson:FR-C311-piece-montee}



\begin{quote}
Before the new one: say the French for the marriage proposal, then the French for to get married.
\end{quote}

\subsection*{You'll want to know: la pièce montée}
\textbf{la pièce montée} — \textquotedblleft{}the tiered wedding cake\textquotedblright{}.

\begin{quote}
À la fin du repas, on a servi la pièce montée.
\end{quote}

\begin{quote}
\emph{At the end of the meal, the wedding cake was served.}
\end{quote}

A tower of small cream puffs, a classic of French weddings.

\begin{cousinweb}
\emph{pièce} is inherited from Late Latin \emph{pettia} \textquotedblleft{}a piece\textquotedblright{}, of Gaulish origin (EtymDB); \emph{montée} is built on \emph{monter} \textquotedblleft{}to go up\textquotedblright{}, from Latin \emph{mōns} \textquotedblleft{}mountain\textquotedblright{} (Lewis \& Short).
\end{cousinweb}

\begin{grammarlens}[title={What you've built}]
One more word for telling, explaining and describing at B1.
\end{grammarlens}

\subsection*{Your turn}
\begin{itemize}
\raggedright
  \item \emph{Say it:} \emph{la pièce montée}
  \item \emph{Say it:} \emph{la pièce montée}, once more
  \item \emph{Say it:} a sentence of your own with \emph{la pièce montée}
  \item \emph{Recall:} say the French for the marriage proposal, then the French for to get married, then say \emph{la pièce montée} again
\end{itemize}

\begin{tcolorbox}[breakable,colback=teal!4,colframe=teal!35!black,title={Before you move on}]
What does la pièce montée mean? (\textquotedblleft{}The tiered wedding cake\textquotedblright{}.) Say it once more.
\end{tcolorbox}
\section[la dragée]{la dragée — the sugared almond}

\label{lesson:FR-C311-dragee}



\begin{quote}
Before the new one: say the French for to get married, then the French for the tiered wedding cake.
\end{quote}

\subsection*{You'll want to know: la dragée}
\textbf{la dragée} — \textquotedblleft{}the sugared almond\textquotedblright{}.

\begin{quote}
Les invités ont reçu des dragées.
\end{quote}

\begin{quote}
\emph{The guests were given sugared almonds.}
\end{quote}

Given at weddings and christenings in France and Belgium.

\begin{cousinweb}
\emph{dragée} comes from Old French \emph{dragie} (EtymDB), usually traced to Latin \emph{tragēmata} \textquotedblleft{}sweets, dessert\textquotedblright{}, a borrowing from Greek.
\end{cousinweb}

\begin{grammarlens}[title={What you've built}]
One more word for telling, explaining and describing at B1.
\end{grammarlens}

\subsection*{Your turn}
\begin{itemize}
\raggedright
  \item \emph{Say it:} \emph{la dragée}
  \item \emph{Say it:} \emph{la dragée}, once more
  \item \emph{Say it:} a sentence of your own with \emph{la dragée}
  \item \emph{Recall:} say the French for to get married, then the French for the tiered wedding cake, then say \emph{la dragée} again
\end{itemize}

\begin{tcolorbox}[breakable,colback=teal!4,colframe=teal!35!black,title={Before you move on}]
What does la dragée mean? (\textquotedblleft{}The sugared almond\textquotedblright{}.) Say it once more.
\end{tcolorbox}
